% ============================================================
%  Chapitre 6 : Modeling — Système RAG pour la documentation SAP
% ============================================================

\chapter{Modélisation}
\section*{Introduction}
% ─────────────────────────────────────────────────────────────
La phase de modélisation constitue l’étape clé de la méthodologie CRISP-DM, où les données préparées sont utilisées pour créer des solutions répondant aux objectifs commerciaux. Dans le contexte de ce projet, cette étape implique la mise en place d'un système de Retrieval-Augmented Generation (RAG) spécifiquement destiné au secteur SAP, afin d’offrir des réponses précises, fiables et contextualisées aux questions techniques.

% ─── Required packages (add to your preamble if not already present) ──────────
% \usepackage{xcolor}
% \usepackage{tikz}
% \usepackage{tcolorbox}
% \usepackage{booktabs}
% \usepackage{colortbl}
% \usetikzlibrary{arrows.meta, positioning, fit, backgrounds, shadows}
% \tcbuselibrary{skins, breakable, most}
%
% ─── Color definitions (add to your preamble) ─────────────────────────────────
% \definecolor{offlinecolor}{HTML}{1565C0}
% \definecolor{onlinecolor}{HTML}{2E7D32}
% \definecolor{offlinebg}{HTML}{E3F2FD}
% \definecolor{onlinebg}{HTML}{E8F5E9}
% \definecolor{arrowcolor}{HTML}{5C6BC0}
% \definecolor{stepbg}{HTML}{EDE7F6}
% \definecolor{stepfg}{HTML}{4527A0}
% ──────────────────────────────────────────────────────────────────────────────

\section{Architecture globale du syst\`{e}me RAG}

% ─────────────────────────────────────────────────────────────────────────────

\subsection{Vue d'ensemble}

Le syst\`{e}me \textbf{RAG} (\textit{Retrieval-Augmented Generation}) con\c{c}u s'articule autour de deux
\'{e}tapes majeures et compl\'{e}mentaires~: une \textcolor{offlinecolor}{\textbf{\'{e}tape \textit{offline}}},
d\'{e}di\'{e}e \`{a} l'ingestion et \`{a} l'indexation des donn\'{e}es, et une
\textcolor{onlinecolor}{\textbf{\'{e}tape \textit{online}}}, responsable du traitement dynamique des
requ\^{e}tes utilisateurs. L'\'{e}tape \textit{offline} convertit les documents SAP bruts en une base
vectorielle interrogeable~; l'\'{e}tape \textit{online} g\`{e}re la requ\^{e}te d'un utilisateur en
l'analysant sur le plan s\'{e}mantique, en extrayant les passages les plus pertinents et en produisant
une r\'{e}ponse structur\'{e}e \`{a} l'aide d'un LLM.

% ── Architecture diagram ──────────────────────────────────────────────────────
\vspace{0.8cm}
\begin{center}
\begin{tikzpicture}[
  node distance=0.5cm and 0.5cm,
  box/.style={
    rectangle, rounded corners=4pt,
    minimum width=2.6cm, minimum height=0.8cm,
    font=\small\bfseries, text centered, drop shadow
  },
  arr/.style={-Stealth, thick},
]
  % Offline row
  \node[box, fill=offlinebg, draw=offlinecolor, text=offlinecolor] (docs)  {Documents SAP};
  \node[box, fill=offlinebg, draw=offlinecolor, text=offlinecolor, right=of docs]  (chunk) {Chunking};
  \node[box, fill=offlinebg, draw=offlinecolor, text=offlinecolor, right=of chunk] (embed) {Embeddings};
  \node[box, fill=offlinebg, draw=offlinecolor, text=offlinecolor, right=of embed] (index) {Index Vectoriel};

  % Online row
  \node[box, fill=onlinebg, draw=onlinecolor, text=onlinecolor, below=1.6cm of docs]  (query)  {Requ\^{e}te};
  \node[box, fill=onlinebg, draw=onlinecolor, text=onlinecolor, right=of query]        (qemb)   {Emb. Requ\^{e}te};
  \node[box, fill=onlinebg, draw=onlinecolor, text=onlinecolor, right=of qemb]         (search) {Similarit\'{e}};
  \node[box, fill=onlinebg, draw=onlinecolor, text=onlinecolor, right=of search]       (llm)    {LLM};
  \node[box, fill=onlinebg, draw=onlinecolor, text=onlinecolor, right=of llm]          (resp)   {R\'{e}ponse};

  % Offline arrows
  \draw[arr, offlinecolor] (docs)  -- (chunk);
  \draw[arr, offlinecolor] (chunk) -- (embed);
  \draw[arr, offlinecolor] (embed) -- (index);

  % Online arrows
  \draw[arr, onlinecolor] (query)  -- (qemb);
  \draw[arr, onlinecolor] (qemb)   -- (search);
  \draw[arr, onlinecolor] (search) -- (llm);
  \draw[arr, onlinecolor] (llm)    -- (resp);

  % Cross-lane: index -> search
  \draw[arr, dashed, arrowcolor, thick]
    (index.south) -- ++(0,-0.4) -| (search.north)
    node[midway, above, font=\tiny\itshape, color=arrowcolor] {r\'{e}cup\'{e}ration};

  % Lane backgrounds
  \begin{scope}[on background layer]
    \node[fit=(docs)(index), inner sep=10pt, fill=offlinebg, draw=offlinecolor,
          dashed, rounded corners=7pt, opacity=0.35] (ob) {};
    \node[fit=(query)(resp), inner sep=10pt, fill=onlinebg, draw=onlinecolor,
          dashed, rounded corners=7pt, opacity=0.35] (nb) {};
  \end{scope}

  \node[left=0.1cm of ob, rotate=90, anchor=south,
        font=\footnotesize\bfseries, color=offlinecolor] {OFFLINE};
  \node[left=0.1cm of nb, rotate=90, anchor=south,
        font=\footnotesize\bfseries, color=onlinecolor] {ONLINE};
\end{tikzpicture}
\end{center}
\vspace{0.4cm}

% ─────────────────────────────────────────────────────────────────────────────


% ─────────────────────────────────────────────────────────────────────────────


% ─────────────────────────────────────────────────────────────────────────────

\subsection{Flux de données}
 
Lorsqu'un utilisateur pose une question dans le chatbot SAP, celle-ci traverse plusieurs étapes avant qu'une réponse lui soit renvoyée. Ce flux peut se décomposer en six grandes phases.

\paragraph{1. Compréhension de la question}
La question brute est d'abord analysée par un module NLP (\textit{Natural Language Processing}) qui identifie l'intention de l'utilisateur (cherche-t-il à résoudre un problème, à comprendre un concept, à suivre une procédure ?), extrait les entités SAP présentes (T-codes, numéros de notes), et reformule la question de manière enrichie pour améliorer la recherche.

\paragraph{2. Recherche hybride dans la base documentaire}
La question enrichie est convertie en vecteur numérique, puis soumise simultanément à deux types de recherche complémentaires~: une \textbf{recherche sémantique} (qui compare le sens de la question avec celui des documents) et une \textbf{recherche par mots-clés} (qui repère les documents contenant exactement les termes SAP de la question). Cette double recherche permet de ne manquer ni les documents conceptuellement proches ni ceux qui contiennent des identifiants exacts.

\paragraph{3. Fusion et classement des résultats}
Les résultats des deux recherches sont fusionnés par un algorithme appelé \textbf{RRF} (\textit{Reciprocal Rank Fusion}), qui combine les classements en favorisant les documents bien placés dans les deux listes à la fois. Un second modèle, le \textbf{cross-encoder}, affine ensuite ce classement en lisant conjointement la question et chaque document candidat pour produire un score de pertinence plus précis.

\paragraph{4. Construction de la réponse.}
Les documents les mieux classés sont injectés dans un prompt structuré transmis au modèle de langage (LLM), qui génère une réponse synthétique en citant ses sources.

\paragraph{5. Contrôle anti-hallucination}
Avant d'être renvoyée à l'utilisateur, la réponse est vérifiée automatiquement~: tout T-code ou numéro de note mentionné dans la réponse mais absent des documents sources est détecté et supprimé.

% ─────────────────────────────────────────────────────────────
\section{Pipeline RAG: Description détaillée}
% ─────────────────────────────────────────────────────────────

\subsection{Vue séquentielle du pipeline}

Le pipeline RAG, implémenté principalement dans \texttt{rag\_pipeline.py}, orchestre l'ensemble du traitement d'une requête utilisateur à travers la fonction centrale \texttt{ask()}. Celle-ci coordonne successivement~: la compréhension de la requête, l'enrichissement sémantique, la recherche hybride, la fusion RRF, le reranking par cross-encoder, la construction du prompt et la génération par le LLM, avant de valider la réponse par un contrôle anti-hallucination.

\subsection{Normalisation de la requête}

La première opération appliquée à toute requête entrante est une normalisation SAP-spécifique réalisée par la fonction \texttt{normalize\_query()}. Cette étape corrige les espaces erronés dans les codes de transaction (ex. \texttt{"SM 37"} $\rightarrow$ \texttt{"SM37"}), et rétablit la casse correcte des termes SAP (\texttt{"abap"} $\rightarrow$ \texttt{"ABAP"}, \texttt{"idoc"} $\rightarrow$ \texttt{"IDoc"}). Cette harmonisation est essentielle pour garantir que les entités extraites correspondront exactement aux valeurs indexées dans la base de données.

\subsection{Compréhension sémantique de la requête (NLP)}
La fonction \texttt{understand\_query()} constitue l'un des éléments clés du pipeline RAG.Cette fonction s'appuie sur une méthode totalement \textit{rule-based} et ne fait appel à aucun modèle de langage (LLM). L'analyse est ainsi entièrement déterministe, rapide et stable, avec un temps d'exécution inférieur à une milliseconde.

Sa principale mission est de convertir une requête utilisateur brute en une forme structurée et enrichie, qui peut être directement utilisée par le moteur de recherche hybride du système.

La fonction renvoie un objet qui inclut plusieurs attributs essentiels :
\begin{itemize}
    \item \textbf{intent} : intention détectée parmi \texttt{lookup}, \texttt{how-to}, \texttt{troubleshoot} ou \texttt{explain}.
    
    \item \textbf{clean\_query} : version nettoyée de la requête, sans mots inutiles ni expressions de remplissage (\textit{filler words}).
    
    \item \textbf{expanded\_query} : requête enrichie par des termes SAP associés afin d’améliorer le rappel documentaire.
    
    \item \textbf{tcodes} : liste des codes de transaction SAP détectés dans la requête.
    
    \item \textbf{sap\_notes} : liste des numéros de notes SAP OSS éventuellement identifiés.
    
    \item \textbf{module} : module SAP détecté (\texttt{BASIS}, \texttt{ABAP}, \texttt{MM}, \texttt{SD}, \texttt{FI}, \texttt{SECURITY}, etc.).
    
    \item \textbf{min\_score} : seuil minimal de similarité cosinus utilisé lors de la recherche vectorielle. Ce seuil est adaptatif : il est fixé à \texttt{0.05} lorsqu’un T-code est détecté dans une requête de type \texttt{lookup}, et à \texttt{0.30} pour les autres cas.
\end{itemize}

\paragraph{Détection de l’intention.}
La détection de l’intention repose sur un ensemble de règles linguistiques appliquées par ordre de priorité~: la première règle qui correspond détermine immédiatement la catégorie. Quatre catégories ont été définies en fonction des besoins réels des utilisateurs SAP~:

\begin{itemize}
  \item \textbf{troubleshoot}~: l’utilisateur signale une erreur ou un dysfonctionnement et cherche à le résoudre~;
  \item \textbf{how-to}~: l’utilisateur demande comment effectuer une procédure ou une action~;
  \item \textbf{lookup}~: l’utilisateur cherche la définition ou le rôle d’un élément précis~; cette catégorie est également assignée automatiquement lorsqu’un T-code est présent dans la question, car une question contenant un T-code est presque toujours une demande d’information ciblée~;
  \item \textbf{explain}~: catégorie par défaut pour toute question conceptuelle ne correspondant pas aux trois cas précédents.
\end{itemize}

\paragraph{Identification du module SAP.}
En parallèle de l’intention, le système identifie le module SAP concerné (BASIS, ABAP, SECURITY, etc.) à partir des mots-clés présents dans la question ou du T-code détecté. Cette information est utilisée comme filtre lors de la recherche documentaire, afin de limiter les résultats aux documents du module pertinent et d’améliorer la précision des réponses.

\paragraph{Continuité conversationnelle.}
Dans un échange multi-tours, l’utilisateur peut poser une question de suivi sans répéter le sujet~: par exemple, après avoir demandé des informations sur \texttt{SM37}, il peut enchaîner avec «~Qu’est-ce qu’il fait~?~». Le système résout automatiquement ce type de pronom en réinjectant le dernier T-code mentionné dans l’historique de conversation, maintenant ainsi la cohérence entre les tours.

Cette étape améliore fortement la qualité des embeddings pour les transactions SAP spécialisées. Par exemple, le T-code \texttt{SMQ1} peut être enrichi par des termes tels que :

\begin{itemize}
    \item qRFC outbound queue ;
    \item asynchronous communication ;
    \item queue monitoring ;
    \item outbound scheduler ;
    \item RFC queue administration.
\end{itemize}

Grâce à cet enrichissement sémantique, le moteur de recherche vectorielle peut retrouver des documents pertinents même lorsque les contenus documentaires ne contiennent pas explicitement le T-code recherché.

\subsection{Vectorisation de la requête}
Le modèle BAAI BAAI/bge-large-en-v1.5 transforme la requête améliorée en un vecteur dense, étant chargé de manière paresseuse via la fonction \texttt{get\_model()} dans un cadre de singleton. La bibliothèque SentenceTransformers assure que l'encodage inclut une normalisation L2. Selon les recommandations du modèle BGE pour la recherche asymétrique (requête brève contre passages longs), un préfixe spécifique \texttt{"Represent this sentence for searching relevant passages: "} est toujours ajouté avant la requête. La taille du vecteur produit est de 1024.

Ce modèle a été sélectionné pour garantir la cohérence entre les embeddings créés lors de l'interrogation et ceux qui sont déjà présents dans la base vectorielle, lesquels ont également été générés avec le modèle BAAI/bge-large-en-v1.5. Cette uniformité assure la compatibilité des représentations vectorielles et renforce la pertinence de la recherche sémantique.




\subsection{Recherche hybride parallèle}

La fonction \texttt{\_do\_search()} orchestre trois recherches simultanées via un \texttt{ThreadPoolExecutor} à trois workers~:

\begin{enumerate}[leftmargin=2em]
  \item \textbf{Recherche vectorielle}~: requête SQL utilisant l'opérateur de distance cosinus \texttt{<=>} de pgvector, avec filtrage par \texttt{min\_score}. Les 15 meilleurs candidats (\texttt{FETCH\_K=15}) sont récupérés.
  \item \textbf{BM25 sur expanded\_query}~: recherche plein texte via \texttt{plainto\_tsquery('english', ...)} et \texttt{ts\_rank()} sur la colonne \texttt{text\_search} (tsvector précalculé).
  \item \textbf{BM25 sur clean\_query}~: même recherche mais sur la requête nettoyée, pour capturer les documents correspondant exactement aux termes SAP.
\end{enumerate}

Les résultats BM25 des deux variantes sont fusionnés par score maximum (\texttt{max}) — et non par somme — pour éviter la surpondération artificielle d'un document présent dans les deux listes .

\subsection{Fusion par Reciprocal Rank Fusion (RRF)}

La recherche hybride produit deux listes de résultats indépendantes~: une issue de la recherche vectorielle, l'autre de la recherche BM25. Le problème est de les combiner intelligemment~: un document peut être en position~3 dans la première liste et en position~1 dans la seconde. Quelle liste croire~?

L'algorithme \textbf{RRF} (\textit{Reciprocal Rank Fusion}) résout ce problème en attribuant à chaque document un score basé sur sa \emph{position} dans chaque liste, et non sur son score brut. Plus un document est bien classé dans \emph{plusieurs} listes à la fois, plus son score RRF est élevé. Un document médiocre dans les deux listes ne peut pas compenser en étant excellent dans une seule.

Dans ce projet, la pondération est délibérément asymétrique~: la recherche BM25 est favorisée ($W_{\text{BM25}} = 0.6$) par rapport à la recherche vectorielle ($W_{\text{vec}} = 0.4$). Ce choix est justifié par la nature du corpus SAP~: les T-codes, noms de modules et numéros de notes sont des identifiants exacts que la recherche par mots-clés repère mieux que la similarité sémantique.

La fonction \texttt{rrf\_fuse()} implémente cet algorithme selon la formule suivante~:

\begin{equation}
  \text{score\_RRF}(d) = \frac{W_{\text{vec}}}{k + r_{\text{vec}}(d)} + \frac{W_{\text{BM25}}}{k + r_{\text{BM25}}(d)}
  \label{eq:rrf}
\end{equation}

où $W_{\text{vec}} = 0.4$, $W_{\text{BM25}} = 0.6$, $k = 60$, et $r(d)$ désigne le rang du document $d$ dans chaque liste. Le choix de favoriser BM25 ($W_{\text{BM25}} > W_{\text{vec}}$) est motivé par la nature terminologique du corpus SAP~: les codes de transaction, les noms de modules et les numéros de notes sont des correspondances exactes que BM25 capture mieux que la sémantique vectorielle. Les 10 meilleurs candidats fusionnés (\texttt{RERANK\_K=10}) sont transmis au cross-encoder.


\subsection{Injection CAG (Context-Augmented Generation)}

Pour certaines questions très ciblées, la recherche dynamique n'est pas toujours nécessaire. Lorsqu'un utilisateur demande directement des informations sur un T-code connu (par exemple \texttt{SM37} ou \texttt{ST22}), les documents les plus pertinents sont prévisibles et peuvent être pré-chargés depuis la base de données, plutôt que recalculés à chaque fois.

C'est le rôle du mécanisme \textbf{CAG} (\textit{Context-Augmented Generation})~: pour les requêtes de type \texttt{lookup} portant sur un T-code identifié, des chunks spécifiques à ce T-code sont récupérés en base et injectés directement en tête du contexte, \emph{avant} les résultats de la recherche hybride. Cela garantit que le contenu le plus autoritaire sur ce T-code est toujours visible par le LLM, même si la recherche dynamique l'aurait classé en dehors des cinq premiers résultats. Ces chunks sont mis en cache après le premier accès, de sorte que chaque T-code ne déclenche qu'une seule requête base de données, quelle que soit la fréquence des questions posées.

Après la fusion RRF et avant le reranking, ce mécanisme injecte des chunks pré-chargés pour les requêtes de type \texttt{lookup} portant sur un T-code connu. La fonction \texttt{get\_cag\_context(tcode, conn\_str)} interroge la base de données pour récupérer les chunks les plus ciblés sur ce T-code (triés par nombre croissant de T-codes mentionnés, pour favoriser les chunks les plus spécifiques)~:

\begin{lstlisting}[language=SQL, caption={Requête CAG : chunks les plus ciblés pour un T-code}]
SELECT ... FROM sap_chunks
WHERE %s = ANY(mentioned_tcodes)
  AND is_parent IS NOT TRUE
ORDER BY array_length(mentioned_tcodes, 1) ASC
LIMIT 6;
\end{lstlisting}

Ces chunks sont mis en cache process-lifetime dans le dictionnaire \texttt{\_CAG\_DB\_CACHE} (protégé par un verrou thread-safe), de sorte que chaque T-code ne déclenche qu'un seul accès DB quelle que soit la fréquence des requêtes. Lors de l'injection, les chunks CAG sont placés \textbf{en tête} du contexte (avant les chunks RAG dynamiques), garantissant que le contenu le plus autoritaire sur ce T-code est toujours visible par le LLM~:

\begin{itemize}[leftmargin=2em]
  \item \textbf{CAG\_MAX\_CHUNKS = 3}~: nombre maximum de chunks pré-chargés injectés.
  \item Les chunks RAG dont l'identifiant est déjà présent dans les chunks CAG sont dédupliqués.
  \item Le total des chunks injectés au LLM reste plafonné à \texttt{TOP\_K} (5 par défaut).
\end{itemize}

% ─────────────────────────────────────────────────────────────
\section{Modélisation des données}
% ─────────────────────────────────────────────────────────────

\subsection{Structure de la table \texttt{sap\_chunks}}

La base de données PostgreSQL conserve la documentation SAP sous forme de chunks dans la table \texttt{sap\_chunks}. Chaque enregistrement correspond à un fragment autonome de documentation, accompagné de métadonnées structurées permettant d’effectuer des recherches ciblées et un filtrage précis des résultats.

\begin{table}[H]
\centering
\caption{Colonnes de la table \texttt{sap\_chunks}}
\label{tab:sap_chunks}

\begin{tabularx}{\textwidth}{|l|l|X|}
\hline
\textbf{Colonne} & \textbf{Type} & \textbf{Description} \\
\hline

\texttt{chunk\_id} &
UUID &
Identifiant unique du chunk \\
\hline

\texttt{text} &
TEXT &
Contenu textuel du chunk \\
\hline

\texttt{title} &
TEXT &
Titre de la section ou du document \\
\hline

\texttt{source\_type} &
VARCHAR &
Type de source : \texttt{sap\_help} ou \texttt{technical\_book} \\
\hline

\texttt{source\_book} &
TEXT &
Nom du livre ou de la documentation source \\
\hline

\texttt{module} &
VARCHAR &
Module SAP (BASIS, ABAP, MM, SD, FI, SECURITY) \\
\hline

\texttt{section\_type} &
VARCHAR &
Type de section (chapter, section, subsection) \\
\hline

\texttt{breadcrumb} &
TEXT[] &
Fil d'Ariane hiérarchique \\
\hline

\texttt{mentioned\_tcodes} &
TEXT[] &
T-codes SAP mentionnés dans le chunk \\
\hline

\texttt{mentioned\_sap\_notes} &
TEXT[] &
Numéros des notes SAP mentionnées \\
\hline

\texttt{embedding} &
VECTOR(1024) &
Vecteur d'embedding dense (BGE-large) \\
\hline

\texttt{text\_search} &
TSVECTOR &
Index de recherche plein texte (BM25) \\
\hline

\texttt{is\_parent} &
BOOLEAN &
Indique si le chunk est un chunk parent \\
\hline

\texttt{parent\_chunk\_id} &
UUID &
Référence au chunk parent (si enfant) \\
\hline

\end{tabularx}
\end{table}
\subsection{Hiérarchie des chunks et stratégie Small-to-Big}

Un problème classique dans les systèmes RAG est le suivant~: pour bien retrouver un document, il vaut mieux des \emph{petits} fragments très précis~; mais pour bien répondre, le LLM a besoin de \emph{grands} blocs de texte avec suffisamment de contexte. Ces deux besoins sont contradictoires si l'on n'utilise qu'un seul niveau de découpage.

La stratégie \textbf{Small-to-Big Retrieval} résout ce problème avec une hiérarchie à deux niveaux~:
\begin{itemize}
  \item Les \textbf{chunks enfants} (petits, ciblés) sont utilisés pour la recherche~: leur précision maximise la pertinence des résultats trouvés.
  \item Les \textbf{chunks parents} (plus larges, plus contextuels) sont transmis au LLM pour la génération~: leur taille fournit suffisamment de contexte pour construire une réponse cohérente.
\end{itemize}

Les colonnes \texttt{is\_parent} et \texttt{parent\_chunk\_id} dans la base de données matérialisent cette hiérarchie. Les requêtes de recherche filtrent systématiquement les chunks parents (\texttt{is\_parent IS NOT TRUE}) pour ne travailler qu'avec les chunks feuilles lors du retrieval.

\section{Techniques d'Embeddings}
% ─────────────────────────────────────────────────────────────

\subsection{Modèle retenu~: BAAI/bge-large-en-v1.5}
Le modèle d’embedding retenu dans le cadre de ce projet est Beijing Academy of Artificial Intelligence BAAI/bge-large-en-v1.5. Développé selon une architecture bi-encodeur, ce modèle génère des représentations vectorielles denses de dimension 1024. Les embeddings produits sont normalisés selon la norme L2, permettant ainsi d’exploiter directement le produit scalaire pour calculer la similarité cosinus entre les vecteurs.
Ce choix est motivé par plusieurs facteurs~:

\begin{itemize}[leftmargin=2em]
  \item \textbf{Performance sur les benchmarks de retrieval}~: BGE-large-en-v1.5 figure parmi les meilleurs modèles sur le benchmark MTEB (\textit{Massive Text Embedding Benchmark}), notamment pour les tâches de recherche d'information en anglais.
  \item \textbf{Optimisation asymétrique}~: le modèle supporte l'ajout d'un préfixe spécifique pour les requêtes (\texttt{"Represent this sentence for searching relevant passages: "}), permettant une séparation sémantique entre la représentation d'une requête courte et celle d'un passage long.
  \item \textbf{Exécution locale}~: le modèle peut être chargé localement via \texttt{SentenceTransformers}, sans dépendance à une API externe payante.
  \item \textbf{Richesse sémantique}~: avec 1024 dimensions, le modèle capture des nuances sémantiques fines, essentielles pour distinguer des concepts SAP proches (ex. \texttt{SM37} vs \texttt{SM36}, \texttt{qRFC} vs \texttt{tRFC}).
\end{itemize}

\subsection{Méthode de vectorisation}

L'embedding des chunks à l'indexation est réalisé sans préfixe (représentation du document), tandis que l'embedding des requêtes utilise le préfixe dédié. La fonction \texttt{embed\_query()} encapsule cette logique~:

\begin{lstlisting}[language=Python, caption={Fonction d'embedding de la requête}]
EMBED_MODEL  = "BAAI/bge-large-en-v1.5"
QUERY_PREFIX = "Represent this sentence for searching relevant passages: "

def embed_query(text: str) -> list[float]:
    return get_model().encode(
        QUERY_PREFIX + text,
        normalize_embeddings=True,
        convert_to_numpy=True
    ).tolist()
\end{lstlisting}

Le chargement du modèle est réalisé en singleton paresseux (\textit{lazy singleton}) via \texttt{get\_model()}, évitant ainsi de recharger le modèle en mémoire à chaque requête.

% ─────────────────────────────────────────────────────────────
\section{Base Vectorielle}
% ─────────────────────────────────────────────────────────────

\subsection{PostgreSQL avec l'extension pgvector}

Le système n'utilise pas un vector store dédié (tel que ChromaDB, Pinecone ou Weaviate), mais s'appuie sur \textbf{PostgreSQL} enrichi de l'extension \textbf{pgvector}. Ce choix architectural présente plusieurs avantages décisifs pour le contexte SAP~:

\begin{itemize}[leftmargin=2em]
  \item \textbf{Unification des données}~: texte, métadonnées structurées (module, T-code, source) et vecteurs cohabitent dans la même table, permettant des filtres combinés en une seule requête SQL.
  \item \textbf{Recherche hybride native}~: PostgreSQL supporte nativement les index \texttt{TSVECTOR} (BM25) et les index vectoriels HNSW via pgvector, éliminant le besoin d'orchestrer deux systèmes distincts.
  \item \textbf{ACID et fiabilité}~: les garanties transactionnelles de PostgreSQL assurent la cohérence des données lors des mises à jour de la base de connaissances.
  \end{itemize}

\subsection{Indexation vectorielle HNSW}

L'extension pgvector crée un index HNSW (\textit{Hierarchical Navigable Small World}) sur la colonne \texttt{embedding}, permettant une recherche approximative du plus proche voisin (\textit{Approximate Nearest Neighbor}) en temps sous-linéaire. La recherche vectorielle est exprimée via l'opérateur \texttt{<=>} qui calcule la distance cosinus~:

\begin{lstlisting}[language=SQL, caption={Requête de recherche vectorielle avec pgvector}]
SELECT chunk_id::text, text, title, source_type, source_book,
       module, section_type, breadcrumb,
       mentioned_tcodes, mentioned_sap_notes,
       is_parent, parent_chunk_id::text,
       1 - (embedding <=> '[0.012,...,0.034]'::vector) AS score
FROM sap_chunks
WHERE is_parent IS NOT TRUE
  AND module = 'BASIS'
ORDER BY embedding <=> '[0.012,...,0.034]'::vector
LIMIT 15;
\end{lstlisting}

Le score de similarité cosinus est calculé comme $1 - \text{distance\_cosinus}$, ce qui donne une valeur entre $-1$ et $1$, les chunks les plus proches du vecteur requête obtenant les scores les plus élevés.


% ─────────────────────────────────────────────────────────────
\section{Mécanisme de Retrieval}
% ─────────────────────────────────────────────────────────────

\subsection{Recherche sémantique vectorielle}
La recherche vectorielle représente le premier axe du retrieval hybride. Elle repose sur la mesure de similarité cosinus entre le vecteur de la requête et ceux des chunks indexés. L'opérateur pgvector \texttt{<=>} effectue cette évaluation de manière optimale grâce à l'index HNSW déjà construit. La fonction \texttt{vector\_search()} élimine les résultats dont le score est inférieur au seuil \texttt{min\_score}, qui oscille entre 0.05 (requêtes \textit{lookup} avec T-code) et 0.30 (requêtes \textit{explain} générales). Cette adaptation dynamique du seuil permet d'ajuster la précision du retrieval selon le type de requête : une question concernant un T-code spécifique doit permettre de récupérer tous les chunks faisant mention de ce code, même avec une similarité globale faible.
% ─────────────────────────────────────────────────────────────
\section{Mécanisme de Reranking par Cross-Encoder}
% ─────────────────────────────────────────────────────────────

\clearpage
\subsection{Modèle cross-encoder ms-marco-MiniLM-L-6-v2}

Le reranking constitue la troisième couche de filtrage du pipeline. Contrairement au bi-encodeur BGE qui encode requête et document séparément, le cross-encoder \texttt{cross-encoder/ms-marco-MiniLM-L-6-v2} analyse conjointement la paire (requête, chunk) pour produire un score de pertinence fine. Ce modèle, entraîné sur le dataset MS-MARCO de recherche de passages, produit des scores sur une échelle approximative de $-10$ à $+10$.

La fonction \texttt{rerank()} applique le cross-encoder sur toutes les paires (requête, chunk candidat) en une seule inférence vectorisée :

\begin{lstlisting}[language=Python, caption={Fonction de reranking par cross-encoder}]
def rerank(query, chunks, top_k=TOP_K):
    scores = get_reranker().predict(
        [[query, c["content"]] for c in chunks],
        show_progress_bar=False
    )
    for chunk, score in zip(chunks, scores):
        chunk["rerank_score"] = round(float(score), 4)
    # Filtrage par seuil MIN_RERANK_SCORE = 2.0
    ranked = [c for c in
              sorted(chunks, key=lambda c: c["rerank_score"],
                     reverse=True)
              if c["rerank_score"] >= MIN_RERANK_SCORE]
    # Deduplication : max 3 chunks par source_book
    source_counts = {}
    deduped = []
    for c in ranked:
        src = (c["metadata"].get("source_book")
               or c["metadata"].get("source_type", ""))
        if source_counts.get(src, 0) < MAX_CHUNKS_PER_SOURCE:
            deduped.append(c)
            source_counts[src] = source_counts.get(src, 0) + 1
    return deduped[:top_k]
\end{lstlisting}

\subsection{Seuillage et déduplication}

Le paramètre \texttt{MIN\_RERANK\_SCORE = 2.0} [F2] filtre les chunks dont le score du cross-encoder est inférieur à ce seuil, éliminant les passages hors-sujet. Ce seuil, relevé de 0.5 à 2.0 par rapport à la version précédente, permet un filtrage efficace tout en conservant les chunks borderline potentiellement utiles. Un mécanisme de déduplication limite à \texttt{MAX\_CHUNKS\_PER\_SOURCE = 3} le nombre de chunks provenant d'une même source, évitant qu'un seul ouvrage ne domine le contexte injecté au LLM.

Un filet de sécurité (\textit{safety net}) est implémenté : si le reranker filtre tous les candidats, les meilleurs candidats RRF bruts sont utilisés à la place, garantissant qu'une réponse est toujours générée.

% ─────────────────────────────────────────────────────────────
% ─────────────────────────────────────────────────────────────
\section{Large Language Model}
% ─────────────────────────────────────────────────────────────

\subsection{Modèle et infrastructure d'exécution locale}
Le modèle de langage utilisé pour générer les réponses est \textbf{llama3.2:3b}, déployé localement à travers \textbf{Ollama} (\texttt{http://localhost:11434/api/generate}). Cette architecture \textit{on-premise} offre plusieurs bénéfices dans le cadre d'un système d'information SAP d'entreprise : sécurité des données (aucune demande ne quitte l'infrastructure interne), contrôle des coûts (pas de facturation à l'utilisation), et autonomie par rapport à un fournisseur de cloud. Un modèle alternatif, \texttt{qwen2.5:7b}, est également pris en charge via la variable d'environnement \texttt{OLLAMA\_MODEL\_QWEN}, permettant de choisir entre rapidité (llama3.2:3b) et qualité supérieure (qwen2.5:7b).

\subsection{Mode streaming et intégration API}

Le LLM supporte deux modes de génération : le mode \textit{streaming} (tokens renvoyés au fil de la génération, affiché en temps réel dans le terminal) et le mode \textit{non-streaming} (réponse complète renvoyée en une fois). L'API FastAPI désactive systématiquement le streaming (\texttt{stream=False}) pour renvoyer une réponse JSON complète au frontend Fiori.

% ─────────────────────────────────────────────────────────────
\section{Prompt Engineering}
% ─────────────────────────────────────────────────────────────

\subsection{Structure du prompt système}

Le prompt système (\texttt{SYSTEM\_PROMPT}) définit le rôle du LLM et ses contraintes de \textbf{synthèse}. Il positionne le modèle comme un \textit{expert SAP technical consultant} et lui demande explicitement de \textit{synthétiser} une réponse à partir des blocs d'evidence pré-sélectionnés (déjà filtrés par \texttt{synthesize\_chunks}), en appliquant quatre règles de synthèse~:

\begin{enumerate}[leftmargin=2em]
  \item \textbf{Fusionner} les faits connexes provenant de plusieurs blocs en une réponse cohérente~; ne pas répéter le même fait par source.
  \item \textbf{Convergence}~: si deux blocs s'accordent sur un fait, l'énoncer une seule fois avec double citation~: \texttt{"as described in [1][3]"}.
  \item \textbf{Conflit}~: si deux blocs se contredisent, présenter explicitement les deux versions~: \texttt{"[1] states X, while [2] states Y"}.
  \item Citer les sources \textbf{en ligne} uniquement via la notation \texttt{[N]}~; ne jamais inventer de citations absentes des en-têtes d'evidence.
\end{enumerate}

Des règles strictes (\textit{STRICT RULES}) interdisent explicitement l'invention de T-codes, de titres de livres, d'URLs SAP Help ou d'étapes non présentes dans l'evidence. Le LLM doit déclarer explicitement~: \texttt{"The available documentation does not cover this."} si le contexte est insuffisant. Cette combinaison de \textit{grounding} strict et de \textit{synthesis rules} constitue la principale défense contre les hallucinations tout en améliorant la cohérence des réponses multi-sources.


\subsection{Construction dynamique du prompt }

Le prompt final soumis au LLM est assemblé dynamiquement par la fonction \texttt{build\_prompt()} à partir de trois composants~: le prompt système fixe, les blocs d’evidence récupérés, et l’historique conversationnel. Chaque chunk est présenté au LLM sous forme d’un bloc numéroté avec un en-tête structuré~:

\begin{lstlisting}[language=Python, caption={Exemple de bloc d’evidence injecté dans le prompt}]
[1] ADM325_Software_Logistics -- Setting Up RFC Connections (BASIS > RFC)
The RFC destination SM59 allows configuring connections between SAP
systems. Enter the target hostname, system number, and logon data...
\end{lstlisting}

L’en-tête contient~: le numéro de citation \texttt{[N]}, le nom du livre source, le titre de la section, et le \textit{breadcrumb} hiérarchique (\texttt{BASIS > RFC}). Cette structuration permet au LLM de citer précisément ses sources via la notation \texttt{[1][3]} et d’éviter toute confusion entre documents de même thématique.

\subsection{Pré-synthèse des chunks (\texttt{synthesis\_text})}

Un document récupéré par le moteur de recherche peut contenir dix phrases~: certaines directement utiles à la question posée, d'autres hors-sujet. Si l'on injecte le document entier dans le prompt du LLM, les phrases inutiles consomment de l'espace et peuvent distraire le modèle. La fenêtre de contexte du LLM étant limitée, chaque phrase compte.

La \textbf{pré-synthèse} résout ce problème~: avant d'injecter un document dans le prompt, on filtre ses phrases pour ne garder que celles qui sont réellement utiles à la question posée. Ce filtrage se fait automatiquement par similarité sémantique~: chaque phrase du document est comparée à la question, et seules les phrases dont le score de pertinence dépasse un seuil sont conservées.

Concrètement, la fonction \texttt{synthesize\_chunks()} applique les étapes suivantes~:

\begin{enumerate}[leftmargin=2em]
  \item Les phrases du chunk sont découpées et vectorisées individuellement avec BGE-large.
  \item Chaque phrase est scorée par similarité cosinus avec le vecteur de la requête.
  \item Seules les phrases dépassant un seuil de pertinence sont conservées dans \texttt{synthesis\_text}.
  \item Les chunks quasi-identiques (similarité inter-chunks $>$ 0.92) sont éliminés comme doublons.
\end{enumerate}

\begin{lstlisting}[language=Python, caption={Utilisation de synthesis\_text dans build\_prompt()}]
# Pre-synthesized sentences if available, fallback to raw content
content = c.get("synthesis_text") or c["content"]
content = content[:500]  # Hard-cap to keep total prompt under 2048 tokens
\end{lstlisting}

Cette pré-synthèse réduit le bruit injecté dans le prompt~: le LLM reçoit une sélection de phrases directement utiles plutôt que des blocs de texte bruts pouvant contenir des informations hors-sujet. Elle améliore également la fidélité de la réponse en concentrant l’attention du modèle sur les passages les plus pertinents.

\subsection{Injection de l’historique conversationnel}

Pour supporter les échanges multi-tours, \texttt{build\_prompt()} intègre les derniers tours de la conversation sous forme d’un bloc \texttt{RECENT CONVERSATION HISTORY}~:

\begin{lstlisting}[language=Python, caption={Injection de l’historique conversationnel}]
=== RECENT CONVERSATION HISTORY ===
User: How do I check background jobs in SAP?
Assistant: To monitor background jobs, use transaction SM37...

=== QUESTION ===
Which statuses indicate a failed job?
\end{lstlisting}

Chaque message est tronqué à 150 caractères afin de contenir la taille totale du prompt dans la fenêtre de contexte du LLM (2048 tokens pour llama3.2:3b). Cette troncature préserve la référence contextuelle tout en évitant de saturer le budget de tokens alloué aux chunks d’evidence.

\subsection{Contrôle anti-hallucination post-génération}

Un modèle de langage peut parfois inventer des informations qui semblent plausibles mais sont factuellement incorrectes~: on parle d'\textbf{hallucination}. Dans le contexte SAP, ce risque est particulièrement problématique~: un T-code inventé ou un numéro de note inexistant peut induire un technicien en erreur directement dans son environnement de production.

Pour limiter ce risque, un mécanisme de vérification automatique est appliqué après chaque génération. Son principe est simple~: tout T-code ou numéro de note SAP mentionné dans la réponse du LLM est comparé à ceux présents dans les documents sources récupérés. Si le LLM cite un identifiant qui n'apparaît dans aucun des documents fournis, cela constitue une hallucination probable.

Ce contrôle opère en deux étapes complémentaires~:

\begin{itemize}
  \item \textbf{Étape 1. Détection par \texttt{check\_hallucinations()}.} La fonction compare les T-codes et numéros de notes SAP présents dans la réponse générée avec ceux effectivement contenus dans les chunks récupérés. Tout T-code connu (\texttt{KNOWN\_TCODES}) mentionné dans la réponse sans apparaître dans le contexte est signalé comme hallucination potentielle~:

\[
\text{T-codes halluc.} = \bigl(\text{T-codes}(\text{réponse}) \setminus \text{T-codes}(\text{contexte})\bigr) \cap \texttt{KNOWN\_TCODES}
\]

  \item \textbf{Étape 2. Sanitisation par \texttt{sanitize\_answer()}.} Les phrases de la réponse dont \textit{tous} les T-codes mentionnés sont hallucinés (non présents dans le contexte) sont supprimées de la réponse finale avant renvoi à l’utilisateur. Les phrases mixtes (T-codes halluc. + T-codes légitimes) sont conservées et signalées en \texttt{warnings} dans la réponse API.
\end{itemize}


% ───────────────────────────────────────────────────
\section{Workflow global d’une requête utilisateur}

Le traitement d’une requête utilisateur suit plusieurs étapes depuis l’interface SAP Fiori jusqu’au retour de la réponse finale. Après la soumission de la question dans le chatbot, la requête est envoyée au backend Python via l’endpoint \texttt{POST /query} de l’API FastAPI. Une session identifiée par UUID est ensuite récupérée ou créée avant l’appel de la fonction principale \texttt{ask()}.

La première étape du pipeline consiste en une analyse NLP \textbf{rule-based} avec \texttt{understand\_query()}, permettant d’identifier l’intention, les entités SAP et une reformulation optimisée de la requête, sans appel au LLM. La requête enrichie est ensuite vectorisée avec BGE-large-en-v1.5 puis comparée au \textbf{cache sémantique}. En cas de similarité suffisante, la réponse est immédiatement renvoyée sans exécuter de retrieval ni de génération.

En cas de cache miss, la fonction \texttt{retrieve()} lance en parallèle une recherche vectorielle et deux recherches BM25. Les résultats sont fusionnés par RRF, avec relâchement progressif des filtres si nécessaire. Pour les requêtes contenant un T-code connu, les chunks du \textbf{cache CAG} sont injectés dans le contexte. Les meilleurs résultats sont ensuite rerankés par un cross-encoder ; un \textit{safety net} restaure les meilleurs chunks RRF si tous les candidats sont rejetés.

Les chunks retenus sont synthétisés afin d’extraire les informations essentielles et supprimer les redondances. Le prompt final est construit avec l’historique conversationnel puis envoyé au modèle Ollama pour la génération de la réponse. Après vérification via \texttt{check\_hallucinations()}, la réponse est stockée dans le cache sémantique et retournée sous forme JSON avec les sources, scores de pertinence, avertissements éventuels et identifiant de session. Enfin, la session est sauvegardée de manière asynchrone dans PostgreSQL.% ─────────────────────────────────────────────────────────────
\section{Algorithme général du pipeline}
% ─────────────────────────────────────────────────────────────

\begin{figure}[H]
\centering
\begin{tikzpicture}[
  font=\scriptsize,
  startstop/.style={draw, rectangle, rounded corners=6pt,
                    fill=SAPBlue, text=white,
                    minimum width=5.5cm, minimum height=0.6cm, text centered},
  process/.style={draw, rectangle,
                  fill=SAPLightBlue, draw=SAPBlue,
                  minimum width=5.5cm, minimum height=0.6cm, text centered},
  cache/.style={draw, rectangle,
                fill=SAPGreen!30, draw=SAPGreen,
                minimum width=5.5cm, minimum height=0.6cm, text centered},
  decision/.style={draw, diamond, aspect=3,
                   fill=SAPLightOrange, draw=SAPOrange,
                   minimum width=5.5cm, minimum height=0.6cm, text centered},
  arr/.style={-{Stealth[length=5pt]}, SAPBlue, thick},
  node distance=0.45cm,
]
\node[startstop] (start)   {Requête utilisateur};
\node[process, below=of start]    (norm)    {1. normalize\_query()};
\node[process, below=of norm]     (nlp)     {2. understand\_query() rule-based + enrich()};
\node[process, below=of nlp]      (embed)   {3. embed(clean\_query) — BGE-large L2};
\node[cache,   below=of embed]    (clookup) {Cache sémantique lookup (seuil 0.92)};
\node[decision, below=of clookup] (chit)    {Cache hit ?};
\node[process, below=of chit]     (search)  {4. vector\_search() \textbar\textbar{} bm25\_search() $\times$2};
\node[decision, below=of search]  (enought) {|résultats| < 3 ?};
\node[process, right=1.5cm of enought] (fallback) {Fallback filtres};
\node[process, below=of enought]  (rrf)     {5. rrf\_fuse() — $W_{v}$=0.4 / $W_{b}$=0.6};
\node[cache,   below=of rrf]      (cag)     {6b. CAG inject (intent=lookup + T-code ?)};
\node[process, below=of cag]      (rerank)  {7. cross\_encoder rerank() — seuil 2.0};
\node[decision, below=of rerank]  (empty)   {chunks vides ?};
\node[process, right=1.5cm of empty] (safety) {Safety net RRF};
\node[process, below=of empty]    (synth)   {8. synthesize\_chunks() — dédup + conflits};
\node[process, below=of synth]    (prompt)  {9. build\_prompt() — synthesis\_text + hist.};
\node[process, below=of prompt]   (gen)     {10. generate() — llama3.2:3b, T=0.1};
\node[process, below=of gen]      (halluc)  {11. check\_hallucinations()};
\node[process, below=of halluc]   (save)    {12. session.add\_turn() async};
\node[cache,   below=of save]     (cstore)  {13. SemanticCache.store()};
\node[startstop, below=of cstore] (end)     {Réponse JSON};

\draw[arr] (start)   -- (norm);
\draw[arr] (norm)    -- (nlp);
\draw[arr] (nlp)     -- (embed);
\draw[arr] (embed)   -- (clookup);
\draw[arr] (clookup) -- (chit);
\draw[arr] (chit.east) -- node[above]{\tiny oui} ++(1.8,0) |- (end.east);
\draw[arr] (chit.south) -- node[right]{\tiny non} (search.north);
\draw[arr] (search)  -- (enought);
\draw[arr] (enought.east) -- node[above]{\tiny oui} (fallback.west);
\draw[arr] (fallback.south) |- (rrf.east);
\draw[arr] (enought.south) -- node[right]{\tiny non} (rrf.north);
\draw[arr] (rrf)     -- (cag);
\draw[arr] (cag)     -- (rerank);
\draw[arr] (rerank)  -- (empty);
\draw[arr] (empty.east) -- node[above]{\tiny oui} (safety.west);
\draw[arr] (safety.south) |- (synth.east);
\draw[arr] (empty.south) -- node[right]{\tiny non} (synth.north);
\draw[arr] (synth)   -- (prompt);
\draw[arr] (prompt)  -- (gen);
\draw[arr] (gen)     -- (halluc);
\draw[arr] (halluc)  -- (save);
\draw[arr] (save)    -- (cstore);
\draw[arr] (cstore)  -- (end);
\end{tikzpicture}
\caption[Pipeline RAG SAP V2]{Flowchart du pipeline RAG SAP V2 Fixed (rag\_pipeline\_v2\_fixed.py)}
\label{fig:pipeline_flowchart}
\end{figure}

\subsection{Limites identifiées}

Malgré ses nombreux atouts, le système comporte certaines limitations. La génération locale avec \texttt{llama3.2:3b} reste plus lente et moins performante qu’un grand modèle hébergé dans le cloud, tel que GPT-4o ou Claude Sonnet. De plus, la fenêtre de contexte limitée à 8192 tokens, bien qu’adaptée à l’utilisation de cinq chunks, peut devenir insuffisante pour des requêtes nécessitant un volume documentaire important.

La détection des hallucinations repose essentiellement sur des heuristiques basées sur des expressions régulières identifiant des T-codes et des numéros de notes. Cette approche ne permet pas de détecter toutes les formes d’hallucination, notamment les affirmations factuellement erronées ne contenant aucun identifiant explicite. Par ailleurs, le \textbf{cache sémantique fonctionne uniquement en mémoire} : son contenu est perdu à chaque redémarrage de l’application et il ne peut pas être mutualisé entre plusieurs instances, ce qui réduit son intérêt dans une architecture distribuée ou multi-processus.

L’analyse \textbf{NLP fondée sur des règles}, bien qu’efficace en termes de rapidité et de déterminisme, montre également ses limites face aux requêtes ambiguës, complexes ou multilingues. Les formulations en français ou les tournures inhabituelles peuvent ainsi conduire à une mauvaise classification des intentions utilisateur. Enfin, les performances du mécanisme de retrieval dépendent fortement de la qualité du découpage initial des documents ainsi que de la pertinence des métadonnées associées, rendant indispensable un important travail de préparation et de structuration des données en amont.
% ─────────────────────────────────────────────────────────────
\section*{Conclusion }
% ─────────────────────────────────────────────────────────────

Ce chapitre a présenté la phase de modélisation du système RAG hybride développé pour la documentation SAP . L’architecture combine plusieurs composants complémentaires, notamment le NLP rule-based, le cache sémantique, le retrieval hybride BM25 + vecteurs, le reranking par cross-encoder et la génération via Llama3.2 avec Ollama. Les améliorations intégrées.
